% SOURCE_STATEMENT: "What is a gapped ground-state phase for an infinite quantum spin system, when phases are defined by differentiable paths of sufficiently short-range interactions with a uniform positive spectral gap?"
% Preparation and final self-audit: CONTRIBUTING.md prompts C and D; source/audit artifacts retained in the local physics_statements_2026_10_11 batch.
% !TeX program = lualatex
\documentclass[11pt]{article}
\usepackage[margin=25mm]{geometry}
\usepackage{fontspec}
\setmainfont{Latin Modern Roman}
\usepackage{amsmath,amssymb,mathtools}
\usepackage{unicode-math}
\setmathfont{Latin Modern Math}
\usepackage{xurl,hyperref}
\hypersetup{colorlinks=true,urlcolor=blue}
\setcounter{secnumdepth}{0}
\setlength{\parindent}{0pt}
\setlength{\parskip}{5pt}
\setlength{\emergencystretch}{3em}
\newcommand{\sref}[2]{\hyperlink{source:#1:#2}{[S#2]}}
\newcommand{\cataloguescope}[1]{\paragraph{Recorded target.}#1}
\newcommand{\R}{\mathbb R}
\newcommand{\C}{\mathbb C}
\newcommand{\N}{\mathbb N}
\newcommand{\Z}{\mathbb Z}
\newcommand{\Q}{\mathbb Q}
\newcommand{\F}{\mathbb F}
\newcommand{\eps}{\varepsilon}
\begin{document}
\section{Classification of Gapped Ground-State Phases}\label{statement}

\subsection{Definitions and mathematical statement}

\paragraph{Shared notation.}$\mathbb N=\{1,2,\ldots\}$; the standard number fields have their usual meanings. All additional conventions are stated locally.


\paragraph{Definitions and conventions.} For a lattice $\Gamma$ with finite-dimensional site Hilbert spaces, $\Phi(X)$ is a short-range interaction supported on finite $X\subseteq\Gamma$, and $H_\Lambda=\sum_{X\subseteq\Lambda}\Phi(X)$ is the local Hamiltonian. Infinite-volume ground states are thermodynamic-limit states on the quasi-local observable algebra. In the source, two ground-state sets belong to the same gapped phase when a differentiable path $\Phi_s$, $s\in[0,1]$, connects their interactions, has the required uniform locality bounds, and has a common positive lower bound $\gamma>0$ for the local ground-state spectral gap.

\paragraph{Source scope and qualifications.} The primary source already supplies the displayed path definition and spectral-flow automorphisms linking equivalent ground-state sets. The broader unresolved task is classification and understanding those phases, not simply inventing the definition. Uniform finite-volume gap bounds, ground-state degeneracy, and optional symmetries must remain explicit.

\paragraph{Catalogue mathematical target.}
\paragraph{Catalogue wording retained for traceability.}
What is a gapped ground-state phase for an infinite quantum spin system, when phases are defined by differentiable paths of sufficiently short-range interactions with a uniform positive spectral gap?

\paragraph{Saved literature context.} The supplied catalogue's assessment dated 2026-08-21 records: The question has a rigorous phase definition and substantial one-dimensional classifications, but no general classification of gapped ground-state phases in the stated setting is known. The assessment is retained as dated provenance; it is not a new comprehensive literature or open-status certification.
\subsection{Short English statement}
Classify the equivalence classes of infinite-volume ground states connected by uniformly gapped paths of local interactions.
\subsection{Sources}
\begin{itemize}
\item[\hypertarget{source:statement:1}{S1}] ULAM.AI and the UnsolvedMath Contributors, UnsolvedMath, catalogue entry. Catalogue curation under CC BY 4.0; definitions and English exposition added locally.
\url{https://huggingface.co/datasets/ulamai/UnsolvedMath}.
\item[\hypertarget{source:statement:2}{S2}] Correlations and Interactions for Random Quantum Systems (2011), Oberwolfach Reports, DOI 10.4171/owr/2011/50, pp. 2837--2890; official report source and mathematical context.
\url{https://doi.org/10.4171/owr/2011/50}.
\item[\hypertarget{source:statement:3}{S3}] Official EMS Press PDF of the same Oberwolfach report; PDF pages 31--33.
\url{https://ems.press/content/serial-article-files/46367}.
\item[\hypertarget{source:statement:4}{S4}] automorphic equivalence theorem; reference in the saved dated catalogue literature review.
\url{https://arxiv.org/abs/1102.0842}.
\item[\hypertarget{source:statement:5}{S5}] one-dimensional classification; reference in the saved dated catalogue literature review.
\url{https://arxiv.org/abs/1212.3718}.
\end{itemize}
\label{end:statement}
\end{document}
